\section{Bounded scale and a limiting fiber map}\label{sec:bounded}

It remains to exclude a subsequence on which $W$ is bounded.
We will choose one complex embedding and normalize each coordinate so
that the uniform laws have nonatomic limits.  The variance estimate will
give the same limits for the fiber laws and for the conditional incoming
laws.  The exact fiber equation then forces a family of M\"obius maps
between these limits that cannot exist.

\subsection{A selected embedding and normalized limits}

At an archimedean place $v$, define the nonnegative local scale
\[
 T_v=\E_{\rm grids}\sum_{j=1}^2
       \bigl(\E_\pi L_j+\E_\alpha S_j\bigr).
\]
Let $V_v$ be the sum of the local integrals on the left of
\eqref{eq:variance}, for the fiber family and the retained incoming
family, including both coordinates and grid expectations but omitting
the place weight.  Since $W$ is bounded, Lemma~\ref{lem:variance} gives
\[
 \sum_{v\mid\infty}\sigma_vT_v\le C,\qquad
 \sum_{v\mid\infty}\sigma_vV_v=\eta_n\longrightarrow0.
\]
The archimedean weights sum to one.  The places with $T_v\le2C+1$
have total weight at least $1/2$, so one of them satisfies $V_v\le2\eta_n$.
Choose such a place for each $n$.  This argument requires no lower
bound on an individual place weight.  Henceforth all coordinates and
grids in this section refer to the selected embedding.

Equation~\eqref{eq:depth} and the grid error bound
\eqref{eq:grid-error} imply, for $j=1,2$,
\[
 \E_{\rm grids}\E_\pi|i_j(x,y)-M_j|
 \le \E_{\rm grids}\E_\pi
       \bigl(L_j(x,y)+S_j(x)+S_j(y)+|\Delta_j(x,y)|\bigr)
 =O(1).
\]
Set $m_j=\E_\pi i_j(x,y)$.  For each realized grid,
$|m_j-M_j|\le\E_\pi|i_j-M_j|$, and thus
$\E_\pi|i_j-m_j|=O(1)$.  Normalize the coordinates by
\[
 U_j(x)=e^{m_j}Z_j(x)+b_j,
\]
where the translation $b_j$ will be chosen below.  We have
\begin{equation}\label{eq:bounded-log}
 \E_\pi\bigl|\log|U_j(x)-U_j(y)|\bigr|\le C_1.
\end{equation}
Only distinct pairs occur here.

Write $\alpha_n^j=(U_j)_*\alpha$ for the empirical coordinate laws.
Equation~\eqref{eq:bounded-log} implies
\[
 (\alpha_n^j\otimes\alpha_n^j)\{|z-z'|>R\}
 \le \frac{C_1}{\log R}\qquad(R>1).
\]
Choose a fixed $R_0>1$ making this bound less than $1/2$.
Averaging over a center at an empirical point gives a radius-$R_0$
disk containing at least half the mass.  Choose $b_j$ to move that
center to zero.  If $R>2R_0$, pairs between this disk and the complement
of the radius-$R$ disk have distance at least $R/2$, so
\begin{equation}\label{eq:bounded-tail}
 \alpha_n^j\{|z|>R\}\le \frac{2C_1}{\log(R/2)}.
\end{equation}
For every center $z_0\in\C$ and $0<r<1/2$, a similar use of
\eqref{eq:bounded-log} gives
\begin{equation}\label{eq:bounded-ball}
 \alpha_n^j(B(z_0,r))^2
 \le \frac1n+\frac{C_1}{\log(1/(2r))}.
\end{equation}
Indeed the diagonal mass of independent pairs in the ball is at most
$1/n$, and every remaining pair has distance at most $2r$.

Let $S=\widehat{\C}$ be the Riemann sphere with its chordal metric,
and write $\mathcal P(S)$ for its probability measures.
The space of probability measures on a compact metric space is compact
and metrizable for weak convergence
\cite[Proposition~5.3 and Theorem~4.2]{WeakCompactness}.
We also use the open-set inequality in the Portmanteau theorem
\cite[Theorem~3.2]{WeakCompactness}. Passing to a common subsequence, let
$\alpha_n^j\Rightarrow\beta_j$ on $S$, for $j=1,2$.
Apply Portmanteau to the open spherical sets
$\{|z|>R\}\cup\{\infty\}$ in \eqref{eq:bounded-tail}; then let
$R\to\infty$.  This shows $\beta_j\{\infty\}=0$.
Apply it to open disks in \eqref{eq:bounded-ball}, then let $r\downarrow0$.
Each $\beta_j$ is therefore a nonatomic probability measure on $\C$.

\subsection{The variance estimate determines the limiting laws}

For either family from Lemma~\ref{lem:variance}, put
$\nu_n^j=(U_j)_*\nu$, and write $\mu\varphi=\int\varphi\,d\mu$
for a measure $\mu$ and a continuous test $\varphi$.
We claim that every $\varphi\in C(S)$ satisfies
\begin{equation}\label{eq:test-convergence}
 \E_{\mathcal F}
 \left|\int\varphi\,d\nu_n^j-\int\varphi\,d\beta_j\right|
 \longrightarrow0.
\end{equation}
The family weights may sum to less than one; each member $\nu$ is
still a probability measure.  In both applications,
$\E_{\mathcal F}\nu\le B\alpha$ for a fixed $B$.

Here are the details converting integrated cell variance into
\eqref{eq:test-convergence}.  Push every realized grid through the same
normalization $z\mapsto e^{m_j}z+b_j$.  This shifts its levels by
$-m_j$ and leaves integration over all levels unchanged.
Fix $R>0$ and $0<\epsilon<1$.  At normalized levels in the interval
\[
 I_\epsilon=[-\log\epsilon,-\log\epsilon+1],
\]
the actual square cells have sides between $\epsilon/(2e)$ and
$\epsilon$.  At most $N=N(R,\epsilon)$ squares meet the radius-$R$
disk, uniformly in the grid and the level.
For each such actual cell $B'$, let $C$ be its transformed occurrence:
$C=B'$ unless $B'$ is the giant, in which case $C=P\setminus B'$.
Since both measures have mass one,
\[
 |\nu(B')-\alpha(B')|=|\nu(C)-\alpha(C)|.
\]
Occurrences are retained with multiplicity.
For those with $|C|>n^{1-c}$, Cauchy--Schwarz over the family, grids,
unit level interval, and at most $N$ occurrences gives an integrated
sum of mean absolute discrepancies at most $\sqrt{NV_v}$.
For the others, including an empty giant complement,
\[
 \E_{\mathcal F}|\nu(C)-\alpha(C)|
 \le(B+1)\alpha(C)\le(B+1)n^{-c}.
\]

Approximate $\varphi$ on each of these squares by its value at any
point of the square.  Its oscillation is at most
$\omega_\varphi(\sqrt2\epsilon)$, where $\omega_\varphi$ is a uniform
modulus of continuity in the spherical metric.  The remaining squares
lie outside the radius-$R$ disk, and their mean family mass is at most
$B\alpha_n^j\{|z|>R\}$.  Integrating the resulting test-function estimate
over grids and the unit interval gives
\begin{align}
 \E_{\mathcal F}|\nu_n^j\varphi-\alpha_n^j\varphi|
 \le{}&
 2\omega_\varphi(\sqrt2\epsilon)
 +(B+1)\|\varphi\|_\infty\alpha_n^j\{|z|>R\}\notag\\
 &+\|\varphi\|_\infty
   \bigl(\sqrt{NV_v}+(B+1)Nn^{-c}\bigr).
 \label{eq:bounded-test-bound}
\end{align}
For fixed $R,\epsilon$, let $n\to\infty$; then let $R\to\infty$ and
$\epsilon\downarrow0$.  Equation~\eqref{eq:bounded-tail} and
$\alpha_n^j\Rightarrow\beta_j$ prove \eqref{eq:test-convergence}.
For the incoming laws, restoring all omitted targets changes a mean
test discrepancy by at most $2\|\varphi\|_\infty$ times their
$o(1)$ total $q_Y$-mass.  Thus \eqref{eq:test-convergence} holds
in full $q_Y$-expectation for $\rho_y$ as well.

We have obtained the same nonatomic coordinate limits for typical
fibers and for typical incoming laws.  The first conclusion will
preserve the fiber equation in a limit; the second will allow its
source coordinate to vary independently of its target.

\subsection{The joint limit and its fiber equation}

For a $q$-edge $(x,y)$ in fiber $e$, consider the joint random object
\[
 \bigl(U_1(x),U_2(x),U_1(y),U_2(y),
       (U_1)_*\mu_e,(U_2)_*\mu_e\bigr)
 \ \in\ S^4\times\mathcal P(S)^2.
\]
Take a subsequential weak limit $Q$, and denote its four coordinate
entries by $(a,b,a',b')$.
The last two entries equal $(\beta_1,\beta_2)$ almost surely.
To justify this simultaneously for all tests, choose a countable
dense sequence $(f_r)$ in the unit ball of $C(S)$.  The metric
\[
 d(\mu,\beta)=\sum_{r\ge1}2^{-r}|\mu f_r-\beta f_r|
\]
defines weak convergence.  Equation~\eqref{eq:test-convergence}
on finitely many tests, followed by a bound on the series tail,
gives convergence of this metric to zero in mean for each fiber
projection.  Passing to $Q$ proves the assertion.

The inequality $q\le B_0\pi\le2B_0\alpha\otimes\alpha$ implies that
each limiting same-coordinate pair law is dominated by
$2B_0\beta_j\otimes\beta_j$.  Indeed the inequality passes to the
limit against every nonnegative continuous test on $S^2$.
Since $\beta_j$ is nonatomic and gives infinity mass zero, we have
\[
 a,b,a',b'\in\C,\qquad a\ne a',\quad b\ne b',
 \qquad a'\in\supp\beta_1,\quad b'\in\supp\beta_2
 \quad\text{$Q$-almost surely}.
\]

Moreover, $a$ has law $\beta_1$ independently of the joint pair
$(a',b')$.  For $f\in C(S)$ and $g\in C(S^2)$, conditioning on the
actual target vertex gives
\begin{align*}
 &\left|\E_q f(U_1(x))g(U_1(y),U_2(y))
     -(\beta_1f)\E_{q_Y}g(U_1(y),U_2(y))\right|\\
 &\hspace{15mm}\le
 \|g\|_\infty\E_{q_Y}
       \left|\int f\,d(U_1)_*\rho_y-\beta_1f\right|
 \longrightarrow0
\end{align*}
by \eqref{eq:test-convergence}.  Such products determine the joint
measure on $S\times S^2$, proving this independence.
Keeping both target coordinates together in $g$ is what allows us to
fix the entire target pair in the final obstruction.

On the open domain of finite quadruples with $a\ne a'$ and $b\ne b'$,
put $t'=(a'-a)(b'-b)$.  We next obtain
\begin{equation}\label{eq:mobius}
 \beta_2=\Phi_*\beta_1,\qquad
 \Phi(z)=b+\frac{t'}{z-a}
 \quad(z\in S),
 \qquad Q\text{-almost surely}.
\end{equation}
For each prelimit edge, its two fiber projection measures are related by
this M\"obius formula, with its four coordinates in place of
$(a,b,a',b')$.  Indeed, normalization multiplies the original fiber product by
$e^{m_1+m_2}$, and evaluation at $y$ identifies the new product as
$(U_1(y)-U_1(x))(U_2(y)-U_2(x))$.

To pass the equality to the limit, note that the matrix
\[
 \begin{pmatrix}b&t'-ab\\1&-a\end{pmatrix}
\]
represents $\Phi$ on the sphere and has determinant $-t'\ne0$.
The map is jointly continuous in its parameters and sphere variable,
including at its pole.  Thus its pushforward acts continuously on
probability measures with the weak topology, as is seen by testing
continuous functions and using uniform continuity on compact parameter
neighborhoods.  The equality graph is relatively closed over the
stated open domain.  Its complement within that domain is open in the
ambient joint space and has zero mass in every prelimit law.
Portmanteau gives it zero $Q$-mass.  The domain itself has full
$Q$-mass, proving \eqref{eq:mobius}.

\subsection{A nonatomic obstruction to the fiber maps}

The following elementary observation completes the argument. The
second source coordinate $b$ may depend on the other coordinates.
In the proof it contributes only an additive constant to a logarithmic
distance, so the independence already proved for $a$ is sufficient.

\begin{lemma}\label{lem:bounded-mobius}
Let $\beta_1,\beta_2$ be nonatomic probability measures on $\C$.
There is no probability law on finite quadruples $(a,b,a',b')$ such
that, almost surely,
\[
 a\ne a',\qquad b\ne b',\qquad
 a'\in\supp\beta_1,\qquad b'\in\supp\beta_2,
\]
$a$ has law $\beta_1$ independently of $(a',b')$, and
\[
 \beta_2=\left(z\mapsto
 b+\frac{(a'-a)(b'-b)}{z-a}\right)_*\beta_1
\]
as measures on the Riemann sphere.
\end{lemma}

\begin{proof}
For a finite real random variable $X$ and a continuous nonnegative
function $w$ supported in a compact subinterval of $(0,1)$, define
\[
 J_w(X)=\int_\R w\bigl(\Pr(X\ge u)\bigr)\,du.
\]
This is finite: the tail probability tends to one at $-\infty$ and
zero at $+\infty$, so the integrand vanishes outside a bounded interval.
It is unchanged when a constant is added to $X$.
We will also use the fact
\begin{equation}\label{eq:bounded-collapse}
 X_k\longrightarrow0\text{ in probability}
 \quad\Longrightarrow\quad J_w(X_k)\longrightarrow0.
\end{equation}
Indeed choose $\delta>0$ such that $w$ vanishes outside
$[\delta,1-\delta]$.  For any $\epsilon>0$, eventually
$\Pr(|X_k|\ge\epsilon)<\delta$.
For $u\ge\epsilon$ the upper tail is below $\delta$, and for
$u\le-\epsilon$ it is above $1-\delta$.  Hence
$0\le J_w(X_k)\le2\epsilon\|w\|_\infty$, proving
\eqref{eq:bounded-collapse}.

Choose continuous weights $w_l:[0,1]\to[0,1]$, for integers $l\ge3$,
each supported inside $(0,1)$ and equal to one on
$[1/l,1-1/l]$.  Write $\Phi$ for the map in the statement.
For $z\ne a,a'$,
\[
 \Phi(z)-b'=(b'-b)\frac{a'-z}{z-a},
 \qquad
 -\log|\Phi(z)-b'|
 =\log\frac{|z-a|}{|z-a'|}-\log|b'-b|.
\]
Nonatomicity removes the exceptional points.  The pushforward
identity and translation invariance of $J_{w_l}$ give
\begin{equation}\label{eq:tail-width}
 \int_\R w_l\!\left(
   \beta_1\!\left\{z:\log\frac{|z-a|}{|z-a'|}\ge u\right\}\right)\,du
 =
 \int_\R w_l\!\left(
   \beta_2\{z:-\log|z-b'|\ge u\}\right)\,du .
\end{equation}
All random variables here are finite almost surely, so the integrals
are finite.  They are measurable in their parameters: define the
logarithmic kernels arbitrarily at coincidences, where the measures
vanish, and apply measurability of parameter integrals to the resulting
Borel functions.

The equality in \eqref{eq:tail-width} depends only on $(a,a',b')$.
Intersect the full-measure sets for all $l$.  The independence
assumption and Fubini's theorem imply that, for almost every target
pair $(a',b')$, all these equalities hold for $\beta_1$-almost every
$a$.  Fix such a target pair in the stated supports.

For $Z$ of law $\beta_2$, the variable $-\log|Z-b'|$ is nonconstant.
More explicitly, nonatomicity gives $r_0>0$ for which
$p=\beta_2\{|z-b'|\ge r_0\}>0$, and support membership gives
$0<r<r_0$ for which $q=\beta_2\{|z-b'|<r\}>0$.
For $-\log r_0<u<-\log r$ its upper tail lies in $[q,1-p]$.
Choose one $l$ with $1/l\le\min(p,q)$.
The right side of \eqref{eq:tail-width} is then strictly positive.

The full-$\beta_1$-measure equality set contains a sequence
$a_k\ne a'$ tending to $a'$, since $a'\in\supp\beta_1$ and
$\beta_1\{a'\}=0$.  For $z$ outside the null countable set
$\{a',a_1,a_2,\ldots\}$,
\[
 \log\frac{|z-a_k|}{|z-a'|}\longrightarrow0.
\]
Thus these variables converge to zero in $\beta_1$-probability.
By \eqref{eq:bounded-collapse}, the left side of
\eqref{eq:tail-width}, for the fixed $l$ just chosen, tends to zero.
The right side is a fixed positive number, a contradiction.
\end{proof}

The law $Q$ constructed above satisfies every hypothesis of
Lemma~\ref{lem:bounded-mobius}, so a bounded-$W$ subsequence is impossible.
Section~\ref{sec:unbounded} ruled out a subsequence with $W\to\infty$.
Every sequence of nonnegative real numbers has a bounded subsequence
or a subsequence tending to infinity.  Therefore the graph sequence
from Proposition~\ref{prop:graph} cannot exist.  The assumption that
$\Theta(s)>0$ somewhere was false, so $F_n(s)\to0$ for each fixed
$s>0$.  This proves Theorem~\ref{thm:main}.
